\bigskip

% course 5, page 1
% \section*{Curs 5 - Pagina 1}

\begin{center}
\fbox{\textbf{ALGEBRA}}\\[2pt]
\fbox{(Curs 5)}
\end{center}
\begin{flushright}
\textit{1.XI.1995}
% \footnote{data exactă e neclară în original -- ar putea fi şi 1.XI, sau anul 1994.}
\end{flushright}

\bigskip
\noindent\underline{\textbf{T.\ lui Lagrange}}

\bigskip
\noindent Fie $(G,\cdot)$ grup.\ $H\leq G$.

\noindent Dacă $x,y\in G$ \quad $x\equiv_s y \pmod H \overset{\text{def}}{\Longleftrightarrow} x^{-1}y\in H$.

\[
R_H^s = \{(x,y)\in G\times G \mid x^{-1}y\in H\} \subseteq G\times G.
\]

\noindent $R_H^s$ = rel.\ de ech.\ pe $G$, ce e compatibilă la st.\ cu $\cdot$.

\[
xR_H^s y \Rightarrow zxR_H^s zy
\]

\noindent Pt.\ $a\in G$, \ $[a]_{R_H^s} = aH = \{ah \mid h\in H\}$ = clasa de resturi la stânga

\bigskip
\noindent $x\equiv_d y \pmod H \overset{\text{def}}{\Longleftrightarrow} xy^{-1}\in H$.

\[
R_H^d = \{(x,y)\in G\times G \mid xy^{-1}\in H\} = \text{rel.\ de ech.\ compat.\ la dr.\ cu } \cdot.
\]

\[
xR_H^d y \Rightarrow xzR_H^d yz.
\]

\[
[a]_{R_H^d} = Ha = \{ha \mid h\in H\}.
\]

\bigskip
\noindent Deci $T=\{a_i\}_{i\in I}$ = transversală la stânga a lui $H$ în $G$

\noindent = sist.\ complet de reprezentanți

\begin{enumerate}
\item[1)] $i\neq j \Rightarrow a_i \not\equiv_s a_j \pmod H$
\item[2)] $(\forall)\, x\in G$ \ $\exists\, i\in I$ a.î.\ $x\equiv_s a_i \pmod H$
\end{enumerate}

\bigskip
\noindent\underline{Not.} $T^{-1} = \{a_i^{-1}\}_{i\in I}$

\noindent $T$ = transv.\ la stânga $\Rightarrow$ $T^{-1}$ = transv.\ la dreapta.

\[
T \xrightarrow{\ } T^{-1}, \qquad a_i \mapsto a_i^{-1} \text{ e bijecție} \ \Rightarrow \ \text{card}\, T = \text{card}\, T^{-1}
\]

\noindent = indicele lui $H$ în $G \overset{\text{not}}{=} [G:H]$

\noindent ord $G$ = card $G$ = $|G|$.

\bigskip
\noindent\fbox{T.\ Lagrange} \quad Dacă $(G,\cdot)$ = grup finit şi $H$ un subgr.\ al lui $G$,

\noindent atunci:

\[
\boxed{|G| = |H|\cdot [G:H]}
\]

\noindent sau, atunci: \quad $\big($ord $G$ = ord $H\cdot [G:H]\big)$

\bigskip
\noindent\underline{Aplicație:} ord $S_n = n!$

\[
S_n = \{\sigma : \{1,2,\ldots,n\}\to\{1,2,\ldots,n\} \mid \sigma \text{ = bijectivă}\}.
\]

% \noindent [Notăm:]\footnote{fragmentul de legătură de aici e foarte compactat în original.}
\[
\sigma = \begin{pmatrix} 1 & 2 & \cdots & n \\ \sigma(1) & \sigma(2) & \cdots & \sigma(n) \end{pmatrix}
\]


